% ============================================================
%  RAPPORT TECHNIQUE — SRUU
%  Système de Réponse aux Urgences Urbaines
%  Sara Arkoub — UMBB M1 Génie Logiciel — 2025-2026
% ============================================================
\documentclass[12pt, a4paper]{report}

% ── Encodage & langue ────────────────────────────────────────
\usepackage[utf8]{inputenc}
\usepackage[T1]{fontenc}
\usepackage[french]{babel}
\usepackage{lmodern}

% ── Mise en page ─────────────────────────────────────────────
\usepackage[top=2.5cm,bottom=2.5cm,left=3cm,right=2.5cm,headheight=15pt]{geometry}

% ── Couleurs ─────────────────────────────────────────────────
\usepackage[dvipsnames, table]{xcolor}
\definecolor{NavyDark}{HTML}{07111F}
\definecolor{NavyMid}{HTML}{0D2137}
\definecolor{NavyCard}{HTML}{102540}
\definecolor{CyanNeon}{HTML}{00C8E0}
\definecolor{OrangeNeon}{HTML}{FF5C00}
\definecolor{GreenNeon}{HTML}{00DFA2}
\definecolor{PurpleNeon}{HTML}{9B59B6}
\definecolor{GrayText}{HTML}{9DBBD4}
\definecolor{LightBg}{HTML}{F0F6FF}
\definecolor{CodeBg}{HTML}{0A1929}

% ── Graphiques ───────────────────────────────────────────────
\usepackage{graphicx}
\usepackage{tikz}
\usetikzlibrary{shapes.geometric, arrows.meta, positioning,
                decorations.pathreplacing, calc, backgrounds, fit}
\usepackage{pgfplots}
\pgfplotsset{compat=1.18}

% ── Boîtes stylées ───────────────────────────────────────────
\usepackage[most]{tcolorbox}
\tcbuselibrary{listings, skins, breakable}

% ── Code source ──────────────────────────────────────────────
\usepackage{listings}
\lstdefinestyle{javaStyle}{
  language=Java,
  basicstyle=\ttfamily\small\color{white},
  keywordstyle=\color{CyanNeon}\bfseries,
  stringstyle=\color{GreenNeon},
  commentstyle=\color{GrayText}\itshape,
  numberstyle=\tiny\color{GrayText},
  numbers=left, numbersep=8pt,
  backgroundcolor=\color{CodeBg},
  frame=none, breaklines=true,
  tabsize=2, showstringspaces=false,
  morekeywords={String, List, Map, HashMap, ConcurrentHashMap,
                AID, ACLMessage, Behaviour, CyclicBehaviour,
                TickerBehaviour, SequentialBehaviour,
                FSMBehaviour, DFAgentDescription, ServiceDescription}
}
\lstdefinestyle{jsonStyle}{
  basicstyle=\ttfamily\small\color{GreenNeon},
  backgroundcolor=\color{CodeBg},
  frame=none, breaklines=true,
  showstringspaces=false,
}
\lstset{style=javaStyle}

% ── En-têtes & pieds de page ─────────────────────────────────
\usepackage{fancyhdr}
\pagestyle{fancy}
\fancyhf{}
\renewcommand{\headrulewidth}{0.4pt}
\renewcommand{\headrule}{\hbox to\headwidth{\color{CyanNeon}\leaders\hrule height\headrulewidth\hfill}}
\fancyhead[L]{\small\color{GrayText}Système de Réponse aux Urgences Urbaines}
\fancyhead[R]{\small\color{GrayText}SRUU — JADE Multi-Agent}
\fancyfoot[C]{\color{GrayText}\small\thepage}

% ── Sections ─────────────────────────────────────────────────
\usepackage{titlesec}
\titleformat{\chapter}[display]
  {\normalfont\huge\bfseries\color{NavyDark}}
  {\color{CyanNeon}\chaptertitlename~\thechapter}{12pt}{\Huge}
\titleformat{\section}
  {\normalfont\Large\bfseries\color{NavyDark}}
  {\color{CyanNeon}\thesection}{1em}{}
\titleformat{\subsection}
  {\normalfont\large\bfseries\color{NavyDark}}
  {\color{OrangeNeon}\thesubsection}{1em}{}
\titleformat{\subsubsection}
  {\normalfont\normalsize\bfseries\color{NavyDark}}
  {\color{GreenNeon}\thesubsubsection}{1em}{}

% ── Sommaire ─────────────────────────────────────────────────
\usepackage{tocloft}
\renewcommand{\cfttoctitlefont}{\Huge\bfseries\color{NavyDark}}
\renewcommand{\cftchapfont}{\bfseries\color{NavyDark}}
\renewcommand{\cftchappagefont}{\bfseries\color{CyanNeon}}
\renewcommand{\cftsecfont}{\color{NavyDark}}
\renewcommand{\cftsecpagefont}{\color{OrangeNeon}}
\renewcommand{\cftsubsecfont}{\small\color{NavyDark}}
\renewcommand{\cftsubsecpagefont}{\small\color{GreenNeon}}
\setlength{\cftbeforechapskip}{8pt}

% ── Liens ────────────────────────────────────────────────────
\usepackage[
  colorlinks=true,
  linkcolor=CyanNeon,
  urlcolor=OrangeNeon,
  citecolor=GreenNeon,
  pdftitle={Rapport SRUU},
  pdfauthor={Sara Arkoub}
]{hyperref}

% ── Tableaux ─────────────────────────────────────────────────
\usepackage{booktabs}
\usepackage{tabularx}
\usepackage{multirow}
\usepackage{array}

% ── Maths ────────────────────────────────────────────────────
\usepackage{amsmath, amssymb}

% ── Utilitaires ──────────────────────────────────────────────
\usepackage{enumitem}
\usepackage{float}
\usepackage{caption}
\captionsetup{font=small, labelfont={bf,color=CyanNeon}, textfont={color=NavyDark}}

% ── Boîtes personnalisées ────────────────────────────────────
\newtcolorbox{infobox}[2][]{
  enhanced, breakable,
  colback=LightBg, colframe=CyanNeon,
  fonttitle=\bfseries\color{NavyDark},
  title=#2, #1,
  attach boxed title to top left={yshift=-2mm, xshift=6mm},
  boxed title style={colback=CyanNeon, colframe=CyanNeon}
}
\newtcolorbox{alertbox}[2][]{
  enhanced, breakable,
  colback=OrangeNeon!8, colframe=OrangeNeon,
  fonttitle=\bfseries\color{white},
  title=#2, #1,
  attach boxed title to top left={yshift=-2mm, xshift=6mm},
  boxed title style={colback=OrangeNeon, colframe=OrangeNeon}
}
\newtcolorbox{successbox}[2][]{
  enhanced, breakable,
  colback=GreenNeon!8, colframe=GreenNeon,
  fonttitle=\bfseries\color{NavyDark},
  title=#2, #1,
  attach boxed title to top left={yshift=-2mm, xshift=6mm},
  boxed title style={colback=GreenNeon, colframe=GreenNeon}
}
\newtcolorbox{codebox}[1]{
  enhanced, breakable,
  colback=CodeBg, colframe=CyanNeon!60,
  fonttitle=\ttfamily\small\color{CyanNeon},
  title=#1,
}

% ────────────────────────────────────────────────────────────
\begin{document}

% ============================================================
%  PAGE DE GARDE
% ============================================================
\begin{titlepage}
\begin{tikzpicture}[remember picture, overlay]

  %% Fond sombre
  \fill[NavyDark] (current page.south west) rectangle (current page.north east);

  %% Panneau gauche
  \fill[NavyMid]
    (current page.north west) -- ++(8cm,0)
    -- ++(0,-5cm) -- ++(-3cm,-23cm)
    -- (current page.south west) -- cycle;

  %% Ligne verticale neon
  \draw[CyanNeon, line width=3pt]
    ([xshift=8cm]current page.north west)
    -- ([xshift=5cm]current page.south west);

  %% Barres horizontales neon haut
  \fill[CyanNeon]
    ([xshift=-4.5cm,yshift=-1.5cm]current page.north east)
    rectangle ([xshift=0cm,yshift=-1.56cm]current page.north east);
  \fill[OrangeNeon]
    ([xshift=-8cm,yshift=-1.5cm]current page.north east)
    rectangle ([xshift=-4.6cm,yshift=-1.56cm]current page.north east);

  %% Barres horizontales bas gauche
  \fill[GreenNeon]
    ([xshift=0cm,yshift=2.2cm]current page.south west)
    rectangle ([xshift=6cm,yshift=2.26cm]current page.south west);
  \fill[CyanNeon]
    ([xshift=0cm,yshift=1.75cm]current page.south west)
    rectangle ([xshift=3.5cm,yshift=1.81cm]current page.south west);

  %% Carré accent bas droite
  \fill[NavyCard]
    ([xshift=-5cm,yshift=0cm]current page.south east)
    rectangle ([xshift=0cm,yshift=4.5cm]current page.south east);
  \draw[PurpleNeon, line width=2pt]
    ([xshift=-5cm,yshift=4.5cm]current page.south east)
    -- ([xshift=0cm,yshift=4.5cm]current page.south east);

  %% Motif losanges haut droite
  \foreach \xi in {0,0.65,1.3}{
    \foreach \yi in {0,0.65,1.3}{
      \fill[NavyCard, opacity=0.7]
        ([xshift=-4.5cm-\xi cm, yshift=-2.2cm-\yi cm]current page.north east)
        -- ++(-0.28,0.28) -- ++(-0.28,-0.28) -- ++(0.28,-0.28) -- cycle;
    }
  }

  %% Institution bas droite
  \node[anchor=south east, xshift=-0.5cm, yshift=0.3cm]
    at (current page.south east)
    {\color{GrayText}\small\ttfamily UMBB -- Département Informatique};

  %% Badge projet haut gauche
  \fill[CyanNeon]
    ([xshift=1cm, yshift=-2.8cm]current page.north west)
    rectangle ([xshift=6.2cm, yshift=-3.35cm]current page.north west);
  \node[anchor=west, xshift=1.2cm, yshift=-3.05cm]
    at (current page.north west)
    {\color{NavyDark}\bfseries\small PROJET SMA -- 2025/2026};

  %% Bandeau pied de page
  \fill[NavyMid]
    (current page.south west) rectangle ([yshift=1.3cm]current page.south east);
  \node[anchor=center, yshift=0.65cm] at (current page.south)
    {\color{GrayText}\small
     9 Agents Autonomes\quad|\quad
     Protocole Contract Net\quad|\quad
     Zéro Deadlock\quad|\quad
     Rapports JSON};

\end{tikzpicture}

\vspace*{2.8cm}
\hspace*{1cm}%
\begin{minipage}{0.65\textwidth}

  {\color{GrayText}\small\bfseries\ttfamily RAPPORT TECHNIQUE}

  \vspace{0.5cm}
  {\fontsize{68}{70}\selectfont\bfseries\color{white} SRUU}

  \vspace{0.2cm}
  {\color{CyanNeon}\rule{11cm}{2pt}}

  \vspace{0.4cm}
  {\Large\color{white}\bfseries Système de Réponse aux Urgences Urbaines}

  \vspace{0.25cm}
  {\large\color{GrayText}\itshape Urban Emergency Response System}

  \vspace{0.9cm}
  {\color{OrangeNeon}\rule{8.5cm}{1pt}}

  \vspace{0.5cm}
  \begin{tabular}{@{}ll}
    \color{GrayText}\small Auteure     & \color{white}\small\bfseries Sara Arkoub \\[5pt]
    \color{GrayText}\small Formation   & \color{white}\small M1 Génie Logiciel \\[5pt]
    \color{GrayText}\small Institution & \color{white}\small UMBB \\[5pt]
    \color{GrayText}\small Année       & \color{white}\small 2025 -- 2026 \\[5pt]
    \color{GrayText}\small Framework   & \color{CyanNeon}\small\bfseries JADE 4.5 + Java 11 \\
  \end{tabular}

  \vspace{1.5cm}
  \foreach \tech/\col in {%
    {Java 11}/CyanNeon,
    {JADE 4.5}/OrangeNeon,
    {FIPA-ACL}/GreenNeon,
    {Contract Net}/PurpleNeon}{%
    \begin{tikzpicture}[baseline]
      \fill[\col] (0,0) rectangle (2.3cm,0.44cm);
      \node[anchor=west] at (0.1cm,0.22cm)
        {\color{NavyDark}\bfseries\footnotesize \tech};
    \end{tikzpicture}\hspace{5pt}%
  }

\end{minipage}
\end{titlepage}

% ============================================================
%  RÉSUMÉ
% ============================================================
\thispagestyle{empty}
\vspace*{2cm}

\begin{tcolorbox}[
  enhanced, breakable,
  colback=LightBg, colframe=CyanNeon, arc=0pt,
  title={\bfseries\large Résumé},
  fonttitle=\color{NavyDark},
  attach boxed title to top left={yshift=-2mm, xshift=8mm},
  boxed title style={colback=CyanNeon, arc=0pt}
]
Ce rapport présente le \textbf{Système de Réponse aux Urgences Urbaines (SRUU)}, un système
multi-agent (SMA) développé avec le framework \textbf{JADE 4.5} pour la coordination intelligente
des ressources d'urgence dans un environnement urbain simulé.

Le système orchestre \textbf{9 types d'agents autonomes} qui communiquent via le protocole
\textbf{FIPA-ACL} et le \textbf{protocole Contract Net} pour détecter, évaluer et résoudre
des incidents de types variés (incendie, urgence médicale, danger biologique). Le
\textit{Dispatcher} central sélectionne l'unité optimale à l'aide d'une
\textbf{fonction d'utilité multi-critères} pondérant la distance, le type de ressource,
le statut et la sévérité de l'incident.

Les résultats de simulation démontrent une coordination robuste avec un temps de décision
inférieur à \textbf{500\,ms}, un temps de réponse moyen de \textbf{45--60 secondes},
et \textbf{zéro deadlock} sur l'ensemble des scénarios testés.
\end{tcolorbox}

\vspace{1cm}

\begin{tcolorbox}[
  enhanced,
  colback=NavyCard!20, colframe=OrangeNeon, arc=0pt,
  title={\bfseries Mots-clés},
  fonttitle=\color{white},
  attach boxed title to top left={yshift=-2mm, xshift=8mm},
  boxed title style={colback=OrangeNeon, arc=0pt}
]
Systèmes Multi-Agents (SMA) $\bullet$
JADE (Java Agent Development Framework) $\bullet$
FIPA-ACL $\bullet$
Contract Net Protocol $\bullet$
Fonction d'utilité $\bullet$
Automate fini (FSM) $\bullet$
Urgences urbaines $\bullet$
Coordination distribuée
\end{tcolorbox}

\clearpage

% ============================================================
%  TABLE DES MATIÈRES
% ============================================================
\tableofcontents
\clearpage

\listoffigures
\clearpage

% ============================================================
%  CHAPITRE 1 — INTRODUCTION
% ============================================================
\chapter{Introduction}

\section{Contexte général}

Les villes modernes font face à une complexité croissante dans la gestion des urgences.
Accidents de la route, incendies, urgences médicales, fuites de matières dangereuses —
ces événements exigent une réponse rapide, coordonnée et adaptée de la part de multiples
services (pompiers, SAMU, police, unités NRBC). La multiplication de ces incidents
et la rareté des ressources rendent la coordination manuelle insuffisante.

Les \textbf{Systèmes Multi-Agents (SMA)} s'imposent comme une solution naturelle à ce défi.
En modélisant chaque acteur du secours comme un agent autonome doté de perceptions,
d'objectifs et de capacités de communication, les SMA permettent une coordination
décentralisée, réactive et robuste, sans point de défaillance unique.

\begin{infobox}{Qu'est-ce qu'un Système Multi-Agent ?}
Un SMA est un ensemble d'agents autonomes qui perçoivent leur environnement, raisonnent
sur leurs objectifs et interagissent via des protocoles de communication formels pour
atteindre des objectifs individuels ou collectifs. Chaque agent est \emph{réactif},
\emph{proactif} et \emph{social}.
\end{infobox}

\section{Objectifs du projet}

Le projet \textbf{SRUU} a pour objectifs :

\begin{itemize}[leftmargin=*, itemsep=5pt]
  \item[\color{CyanNeon}$\blacktriangleright$]
    Concevoir et implémenter un SMA complet avec \textbf{9 types d'agents} autonomes ;
  \item[\color{OrangeNeon}$\blacktriangleright$]
    Mettre en \oe{}uvre le \textbf{protocole Contract Net} (FIPA) pour la négociation
    et l'attribution optimale des ressources ;
  \item[\color{GreenNeon}$\blacktriangleright$]
    Développer une \textbf{fonction d'utilité multi-critères} pour la sélection de l'unité
    la plus appropriée à chaque incident ;
  \item[\color{PurpleNeon}$\blacktriangleright$]
    Simuler des \textbf{scénarios réalistes} incluant des pannes de ressources et des
    incidents simultanés, en garantissant la robustesse du système.
\end{itemize}

\section{Organisation du rapport}

Ce rapport est structuré comme suit :
\begin{itemize}[noitemsep]
  \item Le \textbf{chapitre~\ref{chap:problematique}} analyse la problématique.
  \item Le \textbf{chapitre~\ref{chap:solution}} présente l'architecture SRUU.
  \item Le \textbf{chapitre~\ref{chap:protocoles}} détaille les protocoles et la fonction
    d'utilité.
  \item Le \textbf{chapitre~\ref{chap:simulation}} présente les scénarios et résultats.
  \item Le \textbf{chapitre~\ref{chap:conclusion}} conclut et ouvre des perspectives.
\end{itemize}

% ============================================================
%  CHAPITRE 2 — PROBLÉMATIQUE
% ============================================================
\chapter{Problématique}
\label{chap:problematique}

\section{La gestion des urgences urbaines}

La gestion des urgences en milieu urbain constitue un problème d'optimisation multi-dimen\-sionnel
en temps réel. Plusieurs défis fondamentaux se posent simultanément.

\subsection{Hétérogénéité des ressources}

Une crise urbaine mobilise des ressources de \emph{natures différentes} : ambulances,
camions de pompiers, unités de police, équipes NRBC. Chaque type possède des capacités
spécifiques, une vitesse de déplacement propre et une disponibilité variable.

\begin{alertbox}{Défi n\textdegree 1 — Hétérogénéité}
Comment coordonner des ressources fondamentalement différentes (vitesses, capacités,
disponibilités) dans un cadre unifié et cohérent ?
\end{alertbox}

\subsection{Dynamicité et temps réel}

Les incidents surviennent de manière imprévisible, à tout moment et en tout lieu.
Les décisions d'attribution doivent être prises en quelques centaines de millisecondes,
sans interrompre le traitement des incidents déjà en cours.

\subsection{Incidents simultanés}

Plusieurs incidents peuvent survenir en même temps, créant une compétition pour les
ressources disponibles. Le système doit gérer cette concurrence sans créer de conflits
d'accès (\emph{race conditions}) ni de situations de blocage (\emph{deadlocks}).

\begin{alertbox}{Défi n\textdegree 2 — Concurrence}
Comment garantir qu'une même unité ne soit pas assignée à deux incidents
simultanément, sans mécanisme centralisé bloquant ?
\end{alertbox}

\subsection{Pannes et défaillances}

Les ressources peuvent être défaillantes en cours d'intervention : un camion à court
d'eau, une ambulance en panne. Le système doit détecter ces défaillances et réaffecter
automatiquement les incidents concernés.

\subsection{Optimalité de l'attribution}

L'attribution d'une unité à un incident n'est pas binaire : plusieurs unités peuvent
être candidates. Il faut choisir la \textbf{meilleure} selon plusieurs critères :

\begin{itemize}[leftmargin=*, noitemsep]
  \item La \textbf{proximité géographique} : minimiser le temps de trajet ;
  \item L'\textbf{adéquation du type} : une ambulance pour un incident médical ;
  \item Le \textbf{statut de disponibilité} : préférer une unité libre ;
  \item La \textbf{sévérité de l'incident} : prioriser les cas critiques.
\end{itemize}

\section{Limites des approches classiques}

\begin{table}[H]
\centering
\caption{Comparaison des approches de coordination}
\rowcolors{2}{LightBg}{white}
\begin{tabularx}{\textwidth}{>{\bfseries}lXX}
\toprule
\rowcolor{NavyCard!15}
\color{CyanNeon}Approche & \color{CyanNeon}Avantages & \color{OrangeNeon}Limites \\
\midrule
Humaine centralisée    & Flexibilité, jugement        & Lente, sujette aux erreurs \\
Règles statiques       & Simple, prévisible           & Non adaptative, fragile \\
Optimisation globale   & Optimalité garantie          & NP-difficile, temps prohibitif \\
\textbf{SMA (SRUU)}    & Décentralisé, robuste        & Complexité de conception \\
\bottomrule
\end{tabularx}
\end{table}

\section{Formulation du problème}

\begin{infobox}{Énoncé formel}
Étant donné un ensemble d'incidents $\mathcal{I}$ survenant dynamiquement, chacun
caractérisé par un type, une localisation et une sévérité, et un ensemble d'unités
$\mathcal{U}$ hétérogènes, trouver pour chaque incident $i_k$ l'unité
$u^* \in \mathcal{U}$ maximisant la qualité de la réponse en un délai
$\delta \leq 500\,\text{ms}$, sans deadlock et avec robustesse aux pannes.
\end{infobox}

% ============================================================
%  CHAPITRE 3 — SOLUTION PROPOSÉE
% ============================================================
\chapter{Solution Proposée — Architecture SRUU}
\label{chap:solution}

\section{Vue d'ensemble de l'architecture}

Le SRUU adopte une \textbf{architecture multi-agent hiérarchique à quatre couches},
implémentée avec \textbf{JADE 4.5.0} (Java Agent Development Framework).

\begin{figure}[H]
\centering
\begin{tikzpicture}[
  layer/.style={draw=#1, fill=#1!10, rounded corners=2pt,
                minimum width=14cm, minimum height=1.2cm,
                font=\bfseries},
  lbl/.style={fill=#1, text=NavyDark, rounded corners=1pt,
              font=\small\bfseries, inner xsep=6pt, inner ysep=3pt},
  arr/.style={-{Stealth[length=6pt]}, thick, color=#1}
]
\node[layer=GreenNeon]  (L1) at (0, 4.5)
  {\color{GreenNeon}[PERCEPTION]\quad Sensor Agents ×3
   \hfill\color{GrayText}\small Détection incidents / grille 100×100};
\node[layer=CyanNeon]   (L2) at (0, 3.0)
  {\color{CyanNeon}[DÉCISION]\quad Dispatcher Agent ×1
   \hfill\color{GrayText}\small Contract Net + Fonction d'utilité};
\node[layer=OrangeNeon] (L3) at (0, 1.5)
  {\color{OrangeNeon}[ACTION]\quad Ambulance×2 \;|\; FireTruck×2 \;|\;
   Police×2 \;|\; Biohazard×1
   \hfill\color{GrayText}\small FSM + déplacement};
\node[layer=PurpleNeon] (L4) at (0, 0)
  {\color{PurpleNeon}[SUPPORT]\quad MedCoordinator \;|\;
   TrafficController \;|\; Logger
   \hfill\color{GrayText}\small Services auxiliaires};

\draw[arr=GreenNeon]  (L1.south) -- node[right,font=\small\color{GreenNeon}]{INFORM}         (L2.north);
\draw[arr=CyanNeon]   (L2.south) -- node[right,font=\small\color{CyanNeon}]{CFP / ACCEPT}    (L3.north);
\draw[arr=OrangeNeon] (L3.south) -- node[right,font=\small\color{OrangeNeon}]{INFORM DONE}   (L4.north);
\end{tikzpicture}
\caption{Architecture en couches du système SRUU}
\label{fig:archi}
\end{figure}

\section{Description des agents}

\subsection{Agent Capteur — SensorAgent}

L'agent capteur simule un dispositif de détection d'incidents urbains, déployé à une
position fixe sur la grille.

\begin{itemize}[leftmargin=*, noitemsep]
  \item \textbf{Probabilité de détection :} $p = 0{,}3$ par cycle de 5 secondes
  \item \textbf{Types générés :} FIRE, MEDICAL, BIOHAZARD
  \item \textbf{Sévérité :} entier entre 1 et 10
  \item \textbf{Communication :} message \texttt{INFORM} au Dispatcher
\end{itemize}

\begin{codebox}{Extrait — SensorAgent.java}
\begin{lstlisting}
addBehaviour(new TickerBehaviour(this, detectionInterval) {
    protected void onTick() {
        if (Math.random() < detectionProbability
                && incidentCounter < MAX_INCIDENTS) {
            EmergencyIncident incident = generateRandomIncident();
            ACLMessage inform = new ACLMessage(ACLMessage.INFORM);
            inform.addReceiver(dispatcherAID);
            inform.setContent(incident.toJSON());
            send(inform);
            incidentCounter++;
        }
    }
});
\end{lstlisting}
\end{codebox}

\subsection{Agent Coordinateur — DispatcherAgent}

C'est le \textbf{c\oe{}ur décisionnel} du système. Il reçoit les alertes, lance un
Call For Proposals, évalue les propositions et sélectionne l'unité optimale.

\begin{table}[H]
\centering
\caption{Comportements JADE du DispatcherAgent}
\rowcolors{2}{LightBg}{white}
\begin{tabularx}{\textwidth}{lXl}
\toprule
\rowcolor{NavyCard!15}
\color{CyanNeon}Comportement & \color{CyanNeon}Rôle & \color{CyanNeon}Type JADE \\
\midrule
IncidentReceiver    & Réception des alertes capteurs     & CyclicBehaviour \\
CFPSender           & Envoi du Call For Proposals        & OneShotBehaviour \\
ProposalEvaluator   & Collecte et évaluation des offres  & CyclicBehaviour \\
AssignmentConfirmer & Envoi ACCEPT / REJECT              & OneShotBehaviour \\
AbortHandler        & Gestion des pannes d'unités        & CyclicBehaviour \\
\bottomrule
\end{tabularx}
\end{table}

\subsection{Agents d'Intervention}

\begin{table}[H]
\centering
\caption{Agents d'intervention du système SRUU}
\rowcolors{2}{LightBg}{white}
\begin{tabularx}{\textwidth}{llcXl}
\toprule
\rowcolor{NavyCard!15}
\color{CyanNeon}Agent & \color{CyanNeon}Nb & \color{CyanNeon}Vitesse
  & \color{CyanNeon}Spécialité & \color{CyanNeon}Ressource \\
\midrule
AmbulanceAgent          & 2 & 2\,u/tick & MEDICAL              & --- \\
FireTruckAgent          & 2 & 3\,u/tick & FIRE, BIOHAZARD      & Eau (100\,L) \\
PoliceUnitAgent         & 2 & 3\,u/tick & PERIMETER            & --- \\
BiohazardContainUnit    & 1 & 2\,u/tick & BIOHAZARD, CRYOGENIC & Cap. (50\,u) \\
\bottomrule
\end{tabularx}
\end{table}

\subsection{Agents de Support}

\begin{description}[leftmargin=2cm, style=nextline, itemsep=4pt]
  \item[\color{GreenNeon}MedicalCoordinatorAgent]
    Gère la capacité des hôpitaux, alloue les lits et coordonne les transferts de patients.

  \item[\color{OrangeNeon}TrafficControllerAgent]
    Optimise les routes d'urgence et communique les corridors prioritaires aux unités.

  \item[\color{PurpleNeon}LoggerAgent]
    Enregistre tous les événements et génère des rapports JSON dans \texttt{logs/}.
\end{description}

\section{Modèle d'état fini (FSM) des unités}

\begin{figure}[H]
\centering
\begin{tikzpicture}[
  st/.style={draw=CyanNeon, fill=NavyCard!40, rounded corners=2pt,
             minimum width=2.8cm, minimum height=0.9cm,
             font=\small\bfseries\color{white}},
  ok/.style={draw=GreenNeon, fill=GreenNeon!15},
  ko/.style={draw=OrangeNeon, fill=OrangeNeon!15},
  arr/.style={-{Stealth[length=5pt]}, thick, color=#1, font=\scriptsize}
]
\node[st,ok]  (IDLE)   at (0,0)     {\color{GreenNeon}IDLE};
\node[st]     (RESP)   at (4,0)     {RESPONDING};
\node[st]     (SCENE)  at (8,0)     {ON\_SCENE};
\node[st,ko]  (ABORT)  at (8,-2.4)  {\color{OrangeNeon}ABORTING};
\node[st,ok]  (RET)    at (4,-2.4)  {\color{GreenNeon}RETURNING};

\draw[arr=CyanNeon]   (IDLE)  -- node[above]{\ttfamily ACCEPT}  (RESP);
\draw[arr=CyanNeon]   (RESP)  -- node[above]{\ttfamily ARRIVED} (SCENE);
\draw[arr=GreenNeon]  (SCENE) -- node[right]{\ttfamily DONE}    (ABORT|-SCENE) -- (RET);
\draw[arr=OrangeNeon] (SCENE) -- node[right]{\ttfamily FAILURE} (ABORT);
\draw[arr=OrangeNeon] (ABORT) -- node[below]{\ttfamily sent}    (RET);
\draw[arr=GreenNeon]  (RET)   to[bend right=20] node[below]{\ttfamily BASE} (IDLE);
\end{tikzpicture}
\caption{Automate fini des agents d'intervention (FSM)}
\label{fig:fsm}
\end{figure}

% ============================================================
%  CHAPITRE 4 — PROTOCOLES
% ============================================================
\chapter{Protocoles de Communication et Fonction d'Utilité}
\label{chap:protocoles}

\section{Le protocole Contract Net (FIPA)}

Le \textbf{Contract Net Protocol} (CNP) est un standard FIPA pour la négociation
distribuée entre agents. Il se déroule en cinq étapes :

\begin{figure}[H]
\centering
\begin{tikzpicture}[
  bx/.style={draw=#1, fill=#1!12, rounded corners=2pt,
             minimum width=2.8cm, minimum height=0.9cm,
             font=\small\bfseries\color{#1}},
  ar/.style={-{Stealth[length=5pt]}, thick, color=#1, font=\scriptsize}
]
\node[bx=GreenNeon]  (S1) at (0,0)    {1.\,INCIDENT\\DÉTECTÉ};
\node[bx=CyanNeon]   (S2) at (3.2,0)  {2.\,CFP\\ENVOYÉ};
\node[bx=OrangeNeon] (S3) at (6.4,0)  {3.\,PROPOSE\\REÇU};
\node[bx=CyanNeon]   (S4) at (9.6,0)  {4.\,ACCEPT\\/ REJECT};
\node[bx=GreenNeon]  (S5) at (6.4,-2.2){5.\,MISSION\\TERMINÉE};

\draw[ar=GreenNeon]  (S1) -- node[above]{\ttfamily INFORM}   (S2);
\draw[ar=CyanNeon]   (S2) -- node[above]{\ttfamily CFP}      (S3);
\draw[ar=OrangeNeon] (S3) -- node[above]{\ttfamily PROPOSE}  (S4);
\draw[ar=CyanNeon]   (S4) -- node[right]{\ttfamily ACCEPT}   (S5);
\draw[ar=GreenNeon]  (S5) to[bend right=25] node[left]{\ttfamily INFORM} (S1);
\end{tikzpicture}
\caption{Flux du protocole Contract Net dans SRUU}
\label{fig:cnp}
\end{figure}

\section{Structure des messages FIPA-ACL}

\begin{codebox}{Message PROPOSE — Unité vers Dispatcher}
\begin{lstlisting}[style=jsonStyle]
{
  "performative": "PROPOSE",
  "sender":       "ambulance-agent-1@JADE",
  "receiver":     "dispatcher-agent@JADE",
  "content": {
    "utilityScore":  0.82,
    "distance":      12.5,
    "unitType":      "AMBULANCE",
    "unitStatus":    "IDLE",
    "position":      { "x": 35, "y": 48 }
  },
  "ontology":       "emergency-response-ontology",
  "language":       "FIPA-SL",
  "conversationId": "INC-20260429-001"
}
\end{lstlisting}
\end{codebox}

\section{Fonction d'utilité multi-critères}

\subsection{Définition}

La sélection de l'unité optimale repose sur une \textbf{fonction d'utilité linéaire
pondérée} :

\begin{equation}
  \boxed{
    U(u,\,i) = w_1 \cdot S_{\text{prox}}(u,i)
             + w_2 \cdot S_{\text{type}}(u,i)
             + w_3 \cdot S_{\text{status}}(u)
             + w_4 \cdot S_{\text{sev}}(i)
  }
  \label{eq:utility}
\end{equation}

\begin{table}[H]
\centering
\caption{Poids de la fonction d'utilité}
\rowcolors{2}{LightBg}{white}
\begin{tabularx}{0.80\textwidth}{llXc}
\toprule
\rowcolor{NavyCard!15}
\color{CyanNeon}Param. & \color{CyanNeon}Symbole & \color{CyanNeon}Description & \color{CyanNeon}Poids \\
\midrule
Proximité         & $w_1$ & Distance euclidienne normalisée & $0{,}35$ \\
Type ressource    & $w_2$ & Adéquation unité/incident       & $0{,}25$ \\
Statut            & $w_3$ & IDLE > RETURNING > ACTIVE       & $0{,}20$ \\
Sévérité          & $w_4$ & Priorité de l'incident          & $0{,}20$ \\
\midrule
\multicolumn{3}{r}{\textbf{Total}} & $1{,}00$ \\
\bottomrule
\end{tabularx}
\end{table}

\subsection{Calcul des scores}

\textbf{Score de proximité} (grille $100 \times 100$) :
\[
  S_{\text{prox}}(u,i)
  = 1 - \frac{d(u,i)}{d_{\max}}, \quad
  d(u,i) = \sqrt{(x_u{-}x_i)^2 + (y_u{-}y_i)^2}, \quad
  d_{\max} = 100\sqrt{2}
\]

\textbf{Score de type} :
\[
  S_{\text{type}}(u,i) =
  \begin{cases}
    1{,}0 & \text{si type}(u) \text{ correspond exactement} \\
    0{,}5 & \text{si secondaire admis} \\
    0{,}1 & \text{sinon}
  \end{cases}
\]

\textbf{Score de statut} :
\[
  S_{\text{status}}(u) =
  \begin{cases}
    1{,}0 & \text{IDLE} \\
    0{,}5 & \text{RETURNING} \\
    0{,}2 & \text{ACTIVE}
  \end{cases}
\]

L'unité sélectionnée est $u^* = \arg\max_{u \in \mathcal{U}_{\text{candidats}}} U(u,i)$.

\begin{codebox}{Implémentation — UtilityFunction.java}
\begin{lstlisting}
public static double compute(UnitProposal p, EmergencyIncident i) {
    return DISTANCE_WEIGHT   * computeProximityScore(p, i)
         + TYPE_MATCH_WEIGHT * computeTypeMatchScore(p, i)
         + STATUS_WEIGHT     * computeStatusScore(p.getStatus())
         + SEVERITY_WEIGHT   * (i.getSeverity() / 10.0);
}
\end{lstlisting}
\end{codebox}

\section{Mécanismes de robustesse}

\subsection{Prévention des deadlocks}

Toutes les réceptions de messages utilisent des \texttt{CyclicBehaviour} avec
\texttt{receive()} non-bloquant. Si aucun message n'est disponible, l'agent libère
le processeur (\texttt{block()}) jusqu'à la prochaine activation.

\subsection{Gestion des pannes}

\begin{figure}[H]
\centering
\begin{tikzpicture}[
  bx/.style={draw=CyanNeon, fill=NavyCard!40, rounded corners=2pt,
             minimum width=3.6cm, minimum height=0.75cm, font=\small\color{white}},
  hi/.style={draw=OrangeNeon, fill=OrangeNeon!15, font=\small\bfseries\color{OrangeNeon}},
  ar/.style={-{Stealth[length=5pt]}, thick, color=#1, font=\scriptsize}
]
\node[bx]    (A) at (0,0)    {FireTruck en mission};
\node[bx,hi] (B) at (0,-1.6) {Eau épuisée (0\,L)};
\node[bx]    (C) at (0,-3.2) {ABORT $\to$ Dispatcher};
\node[bx]    (D) at (5.5,-3.2){Dispatcher : nouveau CFP};
\node[bx]    (E) at (5.5,-1.6){FireTruck-2 assigné};
\node[bx]    (F) at (5.5, 0)  {Incident résolu};

\draw[ar=OrangeNeon] (A) -- (B);
\draw[ar=OrangeNeon] (B) -- (C);
\draw[ar=CyanNeon]   (C) -- (D);
\draw[ar=CyanNeon]   (D) -- (E);
\draw[ar=GreenNeon]  (E) -- (F);
\end{tikzpicture}
\caption{Flux de réaffectation après panne de ressource}
\label{fig:abort}
\end{figure}

% ============================================================
%  CHAPITRE 5 — SIMULATION
% ============================================================
\chapter{Simulation et Résultats}
\label{chap:simulation}

\section{Environnement de simulation}

Le système simule un environnement urbain représenté par une \textbf{grille discrète
$100 \times 100$} unités, avec un tick de simulation d'une seconde.

\begin{infobox}{Configuration de la simulation}
\begin{tabular}{@{}ll}
\textbf{Grille :}         & $100 \times 100$ unités \\
\textbf{Tick :}           & 1 seconde \\
\textbf{Capteurs :}       & 3, positions aléatoires \\
\textbf{Incidents max :}  & 3 par capteur (9 total) \\
\textbf{Probabilité :}    & $p = 30\%$ par cycle de 5\,s \\
\textbf{Sévérité :}       & Entier uniforme $[1,\,10]$ \\
\end{tabular}
\end{infobox}

\section{Scénarios testés}

\subsection{Scénario 1 — Incidents multi-types simultanés}

Les trois capteurs déclenchent simultanément : FIRE (sévérité 8), MEDICAL (sévérité 6)
et BIOHAZARD (sévérité 7). Le système assigne automatiquement le bon type d'unité à
chaque incident (matching parfait, score type = 1,0).

\subsection{Scénario 2 — Incidents simultanés en masse}

Les 9 incidents sont générés dans une fenêtre de temps réduite. Le Dispatcher les met
en file d'attente et les unités libérées sont immédiatement réassignées. Aucun deadlock
n'est observé.

\subsection{Scénario 3 — Panne et réaffectation}

\begin{figure}[H]
\centering
\begin{tikzpicture}
\begin{axis}[
  width=11cm, height=5.5cm,
  xlabel={Ticks de simulation},
  ylabel={Eau restante (L)},
  xmin=0, xmax=8, ymin=0, ymax=110,
  xtick={0,1,2,3,4,5,6,7},
  ytick={0,25,50,75,100},
  grid=both, grid style={draw=NavyCard!30},
  axis line style={color=GrayText},
  tick label style={color=GrayText, font=\small},
  label style={color=NavyDark},
  title style={color=NavyDark},
  title={Déplétion de la réserve d'eau — FireTruck-1},
]
\addplot[thick, color=CyanNeon, mark=square*, mark options={fill=CyanNeon}]
  coordinates {(0,100)(1,85)(2,70)(3,55)(4,40)(5,25)(6,10)(7,0)};
\addplot[dashed, thick, color=OrangeNeon]
  coordinates {(7,0)(7,110)};
\node[color=OrangeNeon, font=\small\bfseries] at (axis cs:7,60) [anchor=west] {ABORT};
\end{axis}
\end{tikzpicture}
\caption{Évolution de la réserve d'eau menant à l'abandon de mission}
\label{fig:water}
\end{figure}

Dès l'épuisement de la réserve, FireTruck-1 envoie \texttt{FAILURE} au Dispatcher, qui
lance un nouveau CFP et assigne FireTruck-2 en moins de 2 secondes.

\section{Métriques de performance}

\begin{table}[H]
\centering
\caption{Métriques de performance du système SRUU}
\rowcolors{2}{LightBg}{white}
\begin{tabularx}{0.88\textwidth}{Xcc}
\toprule
\rowcolor{NavyCard!15}
\color{CyanNeon}Métrique & \color{CyanNeon}Mesuré & \color{CyanNeon}Objectif \\
\midrule
Temps de démarrage système     & $< 5$\,s       & $< 10$\,s \\
Temps de décision (Dispatcher) & $100{-}500$\,ms & $< 1$\,s \\
Temps de réponse moyen         & $45{-}60$\,s   & $< 90$\,s \\
Débit                          & $3{-}5$/min    & $> 2$/min \\
Mémoire système                & $\approx 100$\,MB & $< 512$\,MB \\
Deadlocks détectés             & \textbf{0}     & 0 \\
Taux d'attribution réussie     & \textbf{100\%} & $> 95\%$ \\
\bottomrule
\end{tabularx}
\end{table}

\section{Extrait du rapport JSON généré}

\begin{codebox}{logs/incident\_report\_20260429.json}
\begin{lstlisting}[style=jsonStyle]
{
  "generatedAt": "2026-04-29T10:45:00",
  "totalIncidents": 3,
  "statistics": {
    "completed": 3,
    "pending": 0,
    "aborted": 1,
    "averageResponseTimeSeconds": 52.3
  },
  "incidents": [
    {
      "id": "Sensor-1_1714384425123",
      "type": "FIRE",
      "severity": 8,
      "location": { "x": 40, "y": 50 },
      "assignedUnit": "FireTruck-1",
      "status": "COMPLETED",
      "responseTimeSeconds": 45.2
    }
  ]
}
\end{lstlisting}
\end{codebox}

% ============================================================
%  CHAPITRE 6 — CONCLUSION
% ============================================================
\chapter{Conclusion}
\label{chap:conclusion}

\section{Bilan du projet}

Le projet SRUU a atteint l'ensemble de ses objectifs. Nous avons conçu et implémenté
un système multi-agent fonctionnel, capable de coordonner en temps réel une flotte de
ressources hétérogènes pour répondre à des urgences urbaines simulées.

\begin{successbox}{Réalisations clés}
\begin{itemize}[noitemsep]
  \item[\checkmark] 9 agents autonomes fonctionnels avec comportements FSM complets
  \item[\checkmark] Protocole Contract Net FIPA de bout en bout
  \item[\checkmark] Fonction d'utilité multi-critères ($\sum w_i = 1{,}00$)
  \item[\checkmark] Gestion robuste des pannes avec réaffectation automatique
  \item[\checkmark] Zéro deadlock sur tous les scénarios testés
  \item[\checkmark] Rapports JSON complets avec métriques KPI
  \item[\checkmark] Architecture extensible (ajout de nouveaux agents trivial)
\end{itemize}
\end{successbox}

\section{Défis surmontés}

\begin{description}[leftmargin=2cm, style=nextline, itemsep=6pt]
  \item[\color{OrangeNeon}Coordination asynchrone]
    L'utilisation de \texttt{ConcurrentHashMap} et de comportements non-bloquants a
    permis d'éviter les conditions de course et les deadlocks.

  \item[\color{CyanNeon}Calibrage de la fonction d'utilité]
    La détermination des poids optimaux a requis plusieurs itérations de simulation
    pour équilibrer proximité et adéquation du type de ressource.

  \item[\color{GreenNeon}Intégration du 9\textsuperscript{e} agent]
    L'ajout de l'unité NRBC a nécessité d'étendre l'ontologie et la logique du
    Dispatcher pour les types BIOHAZARD et CRYOGENIC\_LEAK.
\end{description}

\section{Enseignements}

\begin{enumerate}[leftmargin=*, itemsep=4pt]
  \item Les \textbf{ontologies formelles} garantissent la sémantique partagée entre
    agents --- toute ambiguïté produit des erreurs silencieuses.
  \item Les \textbf{patterns asynchrones} (CyclicBehaviour + \texttt{block()}) offrent
    une meilleure scalabilité que les approches synchrones bloquantes.
  \item La \textbf{journalisation structurée} (JSON) est indispensable pour l'analyse
    post-simulation et le débogage.
  \item La \textbf{séparation des préoccupations} entre couches (perception, décision,
    action, support) simplifie les tests et l'évolution du système.
\end{enumerate}

\section{Perspectives}

\begin{table}[H]
\centering
\caption{Pistes d'évolution du système SRUU}
\rowcolors{2}{LightBg}{white}
\begin{tabularx}{\textwidth}{lXl}
\toprule
\rowcolor{NavyCard!15}
\color{CyanNeon}Extension & \color{CyanNeon}Description & \color{CyanNeon}Impact \\
\midrule
Dashboard Web          & Interface temps réel avec carte interactive   & UX \\
Optimisation ML        & Apprentissage automatique des poids $w_i$     & Performance \\
Multi-Dispatcher       & Coordinateurs hiérarchiques par zone          & Scalabilité \\
Modèle prédictif       & Anticipation des incidents via historique     & Performance \\
Simulation trafic      & Congestion urbaine affectant les vitesses     & Réalisme \\
Coordination inter-zone & Partage de ressources entre villes           & Scalabilité \\
\bottomrule
\end{tabularx}
\end{table}

\section{Mot de fin}

Ce projet illustre la puissance des Systèmes Multi-Agents pour modéliser et résoudre
des problèmes de coordination complexes en environnement dynamique. L'architecture SRUU,
bien que développée dans un contexte académique, pose des bases solides pour un
déploiement dans des applications réelles de gestion des secours urbains.

\begin{infobox}{Pour aller plus loin}
\textit{«~Les systèmes multi-agents ne résolvent pas des problèmes individuels mieux
qu'un agent unique --- ils résolvent des problèmes qu'aucun agent unique ne pourrait
résoudre.~»}
\end{infobox}

% ============================================================
%  BIBLIOGRAPHIE
% ============================================================
\chapter*{Références}
\addcontentsline{toc}{chapter}{Références}

\begin{enumerate}[leftmargin=*, label={[\arabic*]}, itemsep=5pt]
  \item FIPA. \textit{FIPA Contract Net Interaction Protocol Specification},
    Document SC00029H, 2002.
  \item Wooldridge, M. \textit{An Introduction to MultiAgent Systems}, 2e éd.
    Wiley, 2009.
  \item Bellifemine, F., Caire, G., Greenwood, D.
    \textit{Developing Multi-Agent Systems with JADE}.
    Wiley Series in Agent Technology, 2007.
  \item Shoham, Y., Leyton-Brown, K.
    \textit{Multiagent Systems: Algorithmic, Game-Theoretic, and Logical Foundations}.
    Cambridge University Press, 2009.
  \item JADE Documentation. \textit{JADE Programmer's Guide}, v4.5. TILAB, 2020.
  \item Smith, R.G. \textit{The Contract Net Protocol}.
    IEEE Transactions on Computers, C-29(12), pp.~1104--1113, 1980.
\end{enumerate}

% ============================================================
%  ANNEXES
% ============================================================
\appendix
\chapter{Structure du projet Maven}

\begin{codebox}{Arborescence du projet SRUU}
\begin{lstlisting}[style=jsonStyle]
projet_sma/
├── pom.xml
└── src/main/java/fr/umbb/sma/
    ├── EmergencySystemLauncher.java
    ├── agents/
    │   ├── BaseAgent.java
    │   ├── DispatcherAgent.java
    │   ├── SensorAgent.java
    │   ├── AmbulanceAgent.java
    │   ├── FireTruckAgent.java
    │   ├── PoliceUnitAgent.java
    │   ├── BiohazardContainmentUnitAgent.java
    │   ├── MedicalCoordinatorAgent.java
    │   ├── TrafficControllerAgent.java
    │   └── LoggerAgent.java
    ├── ontology/
    │   ├── EmergencyOntology.java
    │   ├── EmergencyIncident.java
    │   ├── IncidentType.java (FIRE, MEDICAL, BIOHAZARD, PERIMETER)
    │   ├── UnitStatus.java  (IDLE, RESPONDING, ON_SCENE, ABORTING)
    │   ├── Location.java
    │   └── UnitProposal.java
    └── utils/
        ├── IncidentLogger.java
        └── UtilityFunction.java
\end{lstlisting}
\end{codebox}

\chapter{Commandes de compilation et d'exécution}

\begin{codebox}{Compilation et lancement du système}
\begin{lstlisting}[style=jsonStyle]
# Compilation complete
mvn clean package

# Execution standard
java -jar target/sruu-emergency-system-1.0.0.jar

# Execution verbose
JADE_VERBOSE=true java -jar target/sruu-emergency-system-1.0.0.jar

# Lecture du rapport JSON genere
cat logs/incident_report_*.json | python -m json.tool

# Compter les incidents traites
grep -c "\"id\"" logs/incident_report_*.json
\end{lstlisting}
\end{codebox}

\end{document}
